\section{Group Structure of $\mathbf{U(\Z/n\Z)}$ \& Primitive Roots}

Consider the group of units of $\Z/5\Z$, $U(\Z/5\Z)$.

Notice the following exponentiation table for $U(\Z/5\Z)$:

% \begin{center}
%     \begin{tabular}{c|ccccc}
%         $\cdot$ & $1$ & $2$ & $3$ & $4$ \\
%         \hline
%         $1$ & $1$ & $2$ & $3$ & $4$ \\
%         $2$ & $2$ & $4$ & $1$ & $3$ \\
%         $3$ & $3$ & $1$ & $4$ & $2$ \\
%         $4$ & $4$ & $3$ & $2$ & $1$ \\
%     \end{tabular}
% \end{center}
\begin{center}
    \begin{tabular}{c|ccccc}
        $a$ & $a^1$ & $a^2$ & $a^3$ & $a^4$\\
        \hline
        $2$ & $2$ & $4$ & $3$ & $1$\\
        \hline
        $3$ & $3$ & $4$ & $2$ & $1$\\
    \end{tabular}
\end{center}

We can see elements of $U(\Z/5\Z)$ can be generated by $2$ or $3$ by raising powers to them, i.e., \[U(\Z/5\Z) = (\Z/5\Z)^* = \{2^\alpha: \alpha \in \N\} = \{3^\alpha: \alpha \in \N\}.\] Such a group is called \hyperref[def:cyclic]{cyclic group}.

In the above example, $2$ or $3$ generates the group $U(\Z/5\Z)$, but $4$ does not.

\begin{definition}[Primitve Root]
    An integer $a$ is called a \textit{primitive root} modulo $n$ if $\overline{a}$ generates the group $U(\Z/n\Z)$. Equivalently, an integer $a$ is a \textit{primitive root} modulo $n$ if $\phi(n)$ is the smallest positive integer such that $a^{\phi(n)} \equiv 1 \pmod{n}$.
\end{definition}

One can explicitly check that $U(\Z/5\Z), U(\Z/7\Z)$ etc. have primitive roots but $U(\Z/8\Z)$ for example do not have primitive roots.

\begin{theorem}
    If $p$ is prime then $U(\Z/p\Z)$ is a cyclic group.
\end{theorem}
\begin{proof}
    % Let $d\leq p-1$ be the maximal order. Now since the order of an element in $U(\Z/p\Z)$ divides $d$ there are $p-1$ solutions of $x^d=[1]$. So we have $p-1\leq d$, and $d=p-1$.
    To be written.
\end{proof}

The theorem above shows existence of primitive roots modulo all prime numbers. 
